\documentclass[11pt]{article}

\usepackage{amsmath}
\usepackage{amssymb}
\usepackage{amsthm}
\usepackage{geometry}
\usepackage{hyperref}

\geometry{margin=1in}

\title{Theoretical Foundations of GRA-Swarm Systems}
\author{Oleg Bitsoev}
\date{2026}

\newtheorem{definition}{Definition}
\newtheorem{theorem}{Theorem}
\newtheorem{proposition}{Proposition}

\begin{document}

\maketitle

\section{Introduction}

This document provides theoretical foundations for the GRA-Swarm architecture implemented in the GRA-Swarm-Engine repository.

The goal is to formalize the dynamics of resonance-driven multi-agent systems and analyze their collective behavior.

\section{Agent State Space}

Let

\[
S \subset \mathbb{R}^d
\]

be a continuous state space.

Each agent \(i\) has state

\[
s_i \in S
\]

The swarm consists of \(N\) agents:

\[
\mathcal{A} = \{s_1, s_2, ..., s_N\}
\]

\section{Resonance Function}

\begin{definition}[Resonance]

Resonance between agents \(i\) and \(j\) is defined as

\[
R(i,j) = f(s_i, s_j)
\]

where \(f\) is a similarity function.

\end{definition}

In practice we use cosine similarity:

\[
R(i,j) =
\frac{s_i \cdot s_j}{\|s_i\|\|s_j\|}
\]

The resonance matrix is therefore

\[
R \in \mathbb{R}^{N \times N}
\]

\section{Amplification Dynamics}

Each agent receives an amplification score

\[
A_i = \sum_{j=1}^{N} R(i,j)
\]

Agents with larger \(A_i\) are more aligned with the swarm.

This quantity represents the collective support for a given state.

\section{Swarm Evolution Operator}

Let

\[
F : S^N \rightarrow S^N
\]

be the swarm evolution operator.

One iteration of the system consists of:

\begin{enumerate}

\item computing resonance scores \(A_i\)

\item selecting top agents

\item generating mutated agents

\end{enumerate}

Formally,

\[
s_i' = s_{\pi(i)} + \epsilon
\]

where

\begin{itemize}

\item \(\pi(i)\) selects a high scoring agent
\item \(\epsilon\) is stochastic mutation

\end{itemize}

\section{Resonance Graph Interpretation}

The swarm can be interpreted as a weighted graph:

\[
G = (V,E)
\]

where

\begin{itemize}

\item vertices represent agents
\item edges represent resonance strength
\end{itemize}

Edge weight:

\[
w_{ij} = R(i,j)
\]

This graph evolves dynamically as agents mutate.

\section{Spectral Structure}

The resonance matrix \(R\) may exhibit spectral structure.

Let

\[
\lambda_1 \ge \lambda_2 \ge ... \ge \lambda_N
\]

be eigenvalues of \(R\).

Large eigenvalues indicate coherent clusters of agents.

These clusters correspond to regions of consensus in the swarm.

\section{Emergent Clustering}

\begin{proposition}

If the swarm repeatedly selects high resonance agents, clusters of aligned states will form in the state space.

\end{proposition}

\textbf{Sketch:}

Selection favors states with high similarity to others. Mutation introduces local exploration around these states, gradually forming clusters.

\section{Stability Hypothesis}

\begin{theorem}[Informal]

If mutation magnitude is bounded and resonance selection pressure is strong, the swarm converges toward a small number of stable clusters.

\end{theorem}

These clusters represent stable solution families.

\section{Collective Intelligence Hypothesis}

We define collective intelligence in the swarm as the emergence of global search efficiency exceeding that of individual agents.

\begin{definition}

Collective intelligence occurs when

\[
E_{swarm} > E_{individual}
\]

where \(E\) measures problem-solving efficiency.

\end{definition}

This may occur due to:

\begin{itemize}

\item distributed exploration
\item information sharing via resonance
\item evolutionary selection

\end{itemize}

\section{Extension to Language Agents}

If agent states represent language outputs, we obtain a reasoning swarm.

Each agent produces a hypothesis:

\[
h_i
\]

Resonance is computed over semantic similarity:

\[
R(i,j) = similarity(h_i,h_j)
\]

The swarm therefore performs:

\begin{itemize}

\item hypothesis generation
\item hypothesis alignment
\item evolutionary refinement
\end{itemize}

This produces a form of collective reasoning.

\section{Large Scale Swarms}

For large \(N\), the swarm approaches a continuous distribution:

\[
\rho(s,t)
\]

The evolution may be approximated by a stochastic dynamical system.

Future work may explore:

\begin{itemize}

\item mean-field approximations
\item swarm thermodynamics
\item phase transitions in swarm intelligence
\end{itemize}

\section{Conclusion}

This document outlined a theoretical model of resonance-driven swarm systems.

The GRA-Swarm architecture can be interpreted as a dynamical system combining:

\begin{itemize}

\item similarity-based interaction
\item evolutionary selection
\item stochastic exploration
\end{itemize}

Such systems may provide a useful framework for studying emergent distributed intelligence.

\end{document}